\section{Desain}\label{sec:Desain} % (fold)
\subsection{Tech Stack}\label{sub:Tech Stack} % (fold)
Aplikasi ini dirancang menggunakan bahasa Go untuk memudahkan
pengimplementasian \textit{Concurrency} (Spesifikasi Bonus), serta
memanfaatkan sintaks bahasa Go yang mudah dipahami dan mudah
menemukan error. Secara garis besar, arsitektur aplikasi terdiri dari
komponen-komponen berikut.
\begin{enumerate}
    \item Bahasa Pemrograman: Go 1.25.5
    \item Lingkungan Pengembangan: Linux
\end{enumerate}
% subsection Tech Stack (end)

\subsection{Struktur Program}\label{sub:Struktur Program} % (fold)
Proyek disusun mengikuti standar Maven dan spesifikasi tugas:
\begin{enumerate}
    \item \texttt{src/}: Kode sumber program.
    \item \texttt{test/}: Masukan yang digunakan untuk \textit{testing}
    \item \texttt{bin/}: \textit{File executable} atau hasil kompilasi.
    \item \texttt{doc/}: Spesifikasi dan laporan tugas kecil dalam PDF
\end{enumerate}
% subsection Struktur Program (end)
Program utama dibagi menjadi dua komponen utama untuk memisahkan
algoritma dengan GUI:
\begin{itemize}
    \item \texttt{main.go} berisi \textit{entry-point} program yang memanggil fungsi-fungsi lain yang digunakan.
    \item \texttt{octree.go} berisi struktur data \textit{Octree} dan pembangunan \textit{bounding box}
    \item \texttt{build.go} berisi pengimplementasian algoritma \textit{Divide and Conquer} untuk membangun struktur \textit{Octree}, serta penggunaan \textit{concurrency}. 
    \item \texttt{parser.go} berisi logika \textit{parsing} untuk file .obj.
    \item \texttt{viewer.go} berisi pengimplementasian fitur bonus untuk 3D \textit{object viewer}.
    \item \texttt{optimizer.go} berisi logika \textit{optimizing} agar \textit{face} yang muncul tidak duplikat dan sisi kubus yang tertutup kubus lain tidak perlu dimasukkan.
\end{itemize}

\subsection{Cara Kompilasi dan Menjalankan Program}\label{sec:Cara
Kompilasi dan Menjalankan Program} % (fold)
Untuk mengompilasi seluruh kode sumber dan mengunduh dependensi
Go, jalankan perintah berikut pada direktori \textit{src/}.
\begin{lstlisting}[language=bash, basicstyle=\ttfamily\small, breaklines=true]
    go mod tidy
\end{lstlisting}

\noindent Untuk melakukan kompilasi (\textit{build}) program menjadi \textit{file executable}, gunakan perintah berikut.
\begin{lstlisting}[language=bash, basicstyle=\ttfamily\small, breaklines=true]
    go build -o bin/voxelizer ./src/*.go
\end{lstlisting}

\noindent Untuk menjalankan \textit{file executable}, gunakan perintah berikut.
\begin{lstlisting}[language=bash, basicstyle=\ttfamily\small, breaklines=true]
    ./bin/voxelizer
\end{lstlisting}

\noindent Atau jika ingin menjalankan program tanpa mengompilasi, gunakan perintah berikut.
\begin{lstlisting}[language=bash, basicstyle=\ttfamily\small, breaklines=true]
    go run src/*.go
\end{lstlisting}

% subsection Cara Kompilasi dan Menjalankan Program (end)
% section Desain (end)
